\Artigo*%
{filosofia}%
{Érase unha vez un gato...}%
{Mauro Garrido Rodríguez}%
{Como se probaría que está vivo e morto e que significa isto.}%

\begin{multicols}{2}

\subsection*{O infame \textit{gedankenexperiment}}

Todos os estudantes de física do universo coñecido escoitamos algunha vez falar
sobre o paradoxo do gato de Schrödinger e, malia isto, en ningún momento dos
nosos estudos lle demos nin tan sequera un mínimo tratamento. A percepción
xeral, quizais, é que este paradoxo é unha mera popularización da Mecánica
Cuántica ou, peor, que é un problema xa resolto, histórico, pero sen relevancia
actual. O certo é que a través deste \textit{gedankenexperiment} Schrödinger
pretendía sinalar os problemas conceptuais e de conexión co mundo clásico que a
mecánica cuántica (a chamada Interpretación de Copenhague) presentaba e
presenta. Así, nun artigo de $1935$ \cite{schrodinger1} no que discutía a
situación da cuántica estándar e dos seus problemas interpretativos, o
austríaco escribía o seguinte:~\footnote{É de interese resaltar que o paradoxo
non ocupa en absoluto un rol central no seu artigo, é un mero exemplo dentro
dun capítulo onde discute a natureza ontolóxica da función de ondas.}

\Cita[Schöringer]{Un pode incluso establecer casos bastante ridículos. Un gato está
encerrado nunha cámara de aceiro, xunto ao seguinte diabólico dispositivo
(que debe estar protexido fronte á interferencia directa do gato): nun
contador Geiger hai unha pequena cantidade de substancia radioactiva, tan
pequena, que quizais no curso dunha hora un dos átomos decaia, pero tamén,
con igual probabilidade, se cadra ningún; se ocorre, o tubo do contador
descárgase e a través dun relé, libera un martelo que rompe un pequeno frasco
de ácido cianhídrico. Se se deixa que todo este sistema funcione por si só
durante unha hora, diríase que o gato aínda vive se mentres tanto ningún átomo
decaeu. O primeiro decaemento atómico teríao envelenado. A función $\Psi$ do
sistema enteiro expresaría isto contendo o gato vivo e o gato morto
(perdóese a expresión) mesturados ou difusos en partes iguais.}

Así, neste artigo exploraremos o xeito no que podemos trasladar propiedades
cuánticas a obxectos clásicos, se isto é posible e que é o que realmente
queremos dicir cando afirmamos que o gato «está vivo e morto á vez».

\subsection*{Medindo o inmedible}

Como é que un átomo pode «contaxiar» as súas propiedades cuánticas a
un ente macroscópico como un gato? Como sabemos, a desintegración dun átomo
radioactivo é un fenómeno puramente aleatorio: é imposible predicir cando vai
emitir radiación. O átomo, por si só é, polo tanto, intratable como sistema
físico: precisamos facer experimentos con moitos átomos idénticos para obter
resultados dos que extraer estatisticamente conclusións físicas~\footnote{O
mesmo sucederá co gato, será necesario o sacrificio de centos deles en
condicións semellantes de saúde.}. Esta é a gran verdade á que nos enfronta a
mecánica cuántica estándar. Xa, neste intre e sen aínda progresar no
experimento mental, podemos intuír que era o que Schrödinger quería criticar: o
antirrealismo, a noción de que un único átomo é intratable, e que se cadra
extraer conclusións acerca da realidade deste máis alá do seu comportamento
estatístico é algo que pode resultar paradóxico.

\begin{figure}[H]
    \centering
    \includegraphics[width=\linewidth]{revistas/005/imaxes/schrodinger1.png}
    \caption{Alternativas do Gato de Schrödinger. (Foto: Wikipedia)}
    \label{schrodinger1}
\end{figure}

Describamos polo tanto como o átomo inflúe sobre o gato. Supoñamos que o átomo
ten dous estados posibles: excitado e desexcitado. A cuántica ditamina o
seguinte: existe un obxecto matemático, $\Psi$, que codifica, coma un
«catálogo de expectativas» \cite{schrodinger1}, todas as
posibilidades alternativas dun sistema físico concreto, así como a
interferencia entre estas posibilidades. A clave está no de
«interferencia entre posibilidades», posto que senón a cuántica
sería análoga a unha teoría da probabilidade clásica; por exemplo, nada nos
impide facer, para unha moeda
%
\begin{equation}
  \ket{\Psi} = \frac{1}{\sqrt{2}} \ket{\text{cara}} + \frac{1}{\sqrt{2}} \ket{\text{cruz}}.
\end{equation}
%
A diferenza estriba en que as distintas posibilidades do sistema poden dalgún
xeito solaparse, e que este solapamento vén dado por unha fase relativa $\Phi$.
A miúdo facémonos unha imaxe disto co experimento da dobre fenda: o electrón
«pasa» por dous sitios á vez, pensamos nel coma unha onda que se difracta e
interfire consigo mesma tal e como as ondas do mar, incluso o imaxinamos
como esparexido, disolto no espazo. Pero esta imaxe, se ben reconfortante, non
se corresponde coa realidade. O prisma do catálogo de expectativas é máis
axeitado: o que oscila coma unha onda é precisamente a fase de cada
posibilidade; no caso da base de posición, na dobre fenda, a fase asociada a
cada punto do espazo.

\begin{figure}[H]
    \centering
    \includegraphics[width=\linewidth]{revistas/005/imaxes/dobre_fenda.png}
    \caption{Experimento da dobre fenda. (Foto: Wikipedia)}
    \label{dobre_fenda}
\end{figure}

E é máis axeitada porque a interferencia non só se dá no espazo físico, nunha
pantalla que presenta un padrón de interferencia, como en óptica ondulatoria. A
interferencia dáse coas velocidades, co spin, co momento... En calquera base.
E no caso que nos ocupa, entre o estado de excitación dun átomo e a vitalidade
dun gato.

Así, ao principio, o estado do átomo será análogo ao da moeda ($1=$ excitado,
$0=$ desexcitado)
%
\begin{equation}
  \ket{\Psi} = \frac{1}{\sqrt{2}} \ket{1} + \frac{1}{\sqrt{2}} \ket{0}.
\end{equation}
%
Pero algo marca a diferenza: a fase evoluciona co tempo en relación coa
enerxía, \textbf{de acordo coa ecuación de Schrödinger}: é outro xeito de vela,
como a ecuación que rexe a evolución de interferencia entre posibilidades.
Posto que neste caso traballamos precisamente na base de enerxía, cada unha das
posibilidades terá asociada unha fase $\phi = \frac{E}{\hbar}t$, tal que
%
\begin{equation}
  \ket{\Psi} = \frac{1}{\sqrt{2}}( \ket{1} + e^{i\Phi}\ket{0} );~~~\Phi= \frac{E_2 - E_1}{\hbar}t.
\end{equation}
%
Este é o movemento ondulatorio que existe nunha base bidimensional: unha
rotación no plano complexo de frecuencia $\omega = \frac{E}{\hbar}$. Vexamos
como este termo é o responsable de que o gato estea vivo e morto á vez.

Se o átomo se excita, entón o veleno libérase. Se pola contra non o fai, este
non se libera. Isto entrelaza ambas as posibilidades, herdando a fase relativa do
átomo, de xeito que
%
\begin{equation}
  \ket{\Psi} = \frac{1}{\sqrt{2}}( \ket{1, \text{veleno}} + e^{i\Phi}\ket{0, \text{non veleno}} ).
\end{equation}
%
Por último, se o veleno se libera, o gato morre, e se non se libera, entón
mantense vivo: volve suceder o entrelazamento e obtemos~\footnote{Resaltar que
este paradoxo non resolve a discusión de cando é o punto exacto no que morre un
ser vivo: «morto» e «vivo» son etiquetas que designan o estado de cada un
dos átomos do gato cando o veleno é liberado e cando non o é, respectivamente.}
%
\begin{align}
    \ket{\Psi} =& \frac{1}{\sqrt{2}}( \ket{1, \text{veleno}, \text{morto}}\\
               &+ e^{i\Phi}\ket{0, \text{non veleno}, \text{vivo}} ).
\end{align}
%
Pero isto non remata aquí: como accedemos a este estado? A través dun detector.
Podería ser, por exemplo, unha cámara equipada cun software de intelixencia %software: forma máis recomendábel, soporte lóxico
artificial que recoñecese o estado vital do gato (aínda que un experimentador
humano tamén valería). En calquera caso, vemos que sucede o mesmo, que se dá o
entrelazamento
%
\begin{align}
    \ket{\Psi, D} =& \frac{1}{\sqrt{2}}( \ket{1, \text{veleno}, \text{morto}, M} \\
                  &+ e^{i\Phi}\ket{0, \text{non veleno}, \text{vivo}, V} ).
\end{align}

\begin{figure}[H]
    \centering
    \includegraphics[width=\linewidth]{revistas/005/imaxes/schrodinger2.jpg}
    \caption{Detector de estado vital.}
    \label{schrodinger2}
\end{figure}

Por que, entón, obtemos un único resultado, se a cadea de entrelazamentos non
cesa? Este é o chamado \textbf{problema da medida}, que Schrödinger tamén
aborda no artigo citado, e que clasicamente, dentro da Interpretación de
Copenhague, se resolvía a golpe de postulado, o de colapso. Abordaremos este
problema na seguinte sección, pois garda relación coa imposibilidade técnica de
realizar este experimento mental.

Por agora contentémonos co colapso: cando un mide, o estado decídese por unha
posibilidade ou outra. Pero, que diferenza existe entón co caso da moeda
clásica, se sempre que medimos o estado do gato, está vivo \textit{ou} morto? A
fase $\Phi$ está presente, pero non se manifesta nas nosas medidas. Precisamos
un observable que sexa sensible a esta fase. Para simplificar, vexamos primeiro
que instrumento temos que empregar para observar a interferencia do átomo:

\begin{figure}[H]
    \centering
    \includegraphics[width=\linewidth]{revistas/005/imaxes/interferometro.jpg}
    \caption{Porta de Hadamard con detectores de excitación.}
    \label{interferometro}
\end{figure}

Esta é a chamada «porta de Hadamard». No noso caso, consiste nunha cavidade
que, mediante radiación electromagnética, pode facer que o átomo se excite por
absorción ou que se desexcite por emisión estimulada coa mesma probabilidade. A
base na que mide é precisamente a que combina estado excitado con
desexcitado~\footnote{O signo menos no segundo estado de medida é debido a que
ambos deben de ser ortogonais e á súa vez ser unha combinación de excitado e
desexcitado.}
%
\begin{align}
   \ket{1} \rightarrow & \frac{1}{\sqrt{2}} (\ket{0} + \ket{1}), \\
   \ket{0} \rightarrow & \frac{1}{\sqrt{2}} (\ket{0} - \ket{1}).
\end{align}
%
Se sometemos o átomo a esta medida, o estado entrelazado quedará como
%
\begin{equation}
    \ket{\Psi} = \frac{1}{2} \left[ (1+e^{i\Phi})\ket{0} + (1-e^{i\Phi})\ket{1} \right].
\end{equation}
%
De xeito que, se nos fixamos, por exemplo, na probabilidade de que o medidor de
estado excitado se active
%
\begin{equation}
    |\braket{1|\Psi}|^2 = \sin{}^2(\Phi/2).
\end{equation}
%
Obtemos interferencia! A existencia da fase fará que o átomo estea excitado con
máis ou menos probabilidade. O análogo poderiamos facer co gato, sempre e cando
dispuxésemos dun dispositivo que puidese reconfigurar os átomos do gato dun
estado vital a outro, isto é: revivilo.

\subsection*{A interpretación da superposición}

Pero supoñamos que podemos revivir o gato. Que demostrariamos con este
experimento? Demostrariamos que o felino, antes de medilo, non está nin vivo
(pois só con $\ket{\Psi} = \ket{0, \text{non veleno}, \text{vivo}, V\ }$ non
observamos interferencia) nin morto (con $\ket{\Psi} = \ket{1, \text{veleno},
\text{morto}, M\,}$ tampouco podemos observar interferencia) \textit{nin dous
gatos, un vivo e outro morto} (temos a certeza de que sempre empregamos un
gato) nin, por suposto, un gato nin vivo nin morto. O problema é que estas son
todas as opcións lóxicas. O único que podemos facer é postular un novo modo de
\textit{ser} que chamaremos \textbf{superposición} \cite{albert1992}. Este é
realmente o maior misterio ao que nos enfronta a Mecánica Cuántica, que este
experimento mental leva máis alá: \textbf{\textit{Que se sente estar en
superposición?}} É esta a pregunta central das interpretacións da mecánica
cuántica.

\subsection*{Estamos en superposición? A decoherencia}

Como xa adiantamos, o proceso de medida xa non se entende actualmente co
postulado do colapso: este levaba a considerar os aparellos de medida como
algo en esencia diferente, e incluso a postular que era a consciencia humana a
que provocaba os colapsos. A moderna teoría da decoherencia simplemente afirma
o seguinte: todo está en superposición, pois é o que dita a ecuación de
Schrödinger, \textbf{pero a coherencia} (isto é, a fase relativa entre estados)
\textbf{pérdese debido á interacción co ambiente}.

O ambiente actúa coma unha especie de medidor: que o veleno sexa liberado ou
non, por exemplo, causa estados totalmente distintos do ambiente: isto provoca
que a información de se se liberou se filtre ao ambiente, quede aí; nese intre
o entrelazamento esténdese entre todas as moléculas de aire, e é coma se
tivésemos executado unha medida. O gato está ou vivo ou morto \footnote{E xa
nin falemos da decoherencia entre os estados vitais do gato.}. Pero é
importante recalcar o seguinte: todo segue en superposición, só que entre
moitísimos graos de liberdade. Se executásemos unha medida análoga á porta de
Hadamard pero considerando todo o ambiente (o cal é a todas luces
extremadamente difícil), recuperariamos de novo esa fase $\Phi$. Se non
queremos volver postular o colapso, todo nos conduce ao seguinte: unha función
de onda global $\Psi$ para todo o Universo.

\textit{E ben, como se sente estar en superposición?}

\printbibliography

\end{multicols}
